\chapter{Introduction}

% Putting prospectus in here for now

By 2050, over 60\% of the global population will live in cities, making public transit not just essential but foundational to the functioning of modern society. Yet, despite advances in mobility technology, cities still struggle to design bus systems that are efficient, equitable, and adaptable to evolving urban needs. Bus service design aims at identifying bus routes which match historic ridership demand while complying with requirements from the built environment, the expected reliability, equity, and efficiency, as well as the environmental impact.

In 1984, \textcite{ceder_bus_1986} first attempted to apply an algorithmic approach to the design of urban bus networks. Their article highlights how most prior work has been put into optimizing bus scheduling rather than the network itself, and they propose a multi-level approach to solve this. Extending this, \textcite{zardini_co-design_2023} proposes an interdependent solution, where different modes must continuously balance each other through a variety of factors, utilizing category theory to do so. In addition to this, \textcite{furth_setting_1981} delved into the optimization of bus transit to maximize social benefit. Yet, they note the need for a more sophisticated approach, such as utilizing the applied category theory framework. 

Very little work has been done on integrating equity analyses into these computational network optimizations. Furthermore, as demonstrated by \textcite{situ_is_2025}, any equity analysis required by Title VI doesn’t accurately capture the nuances of large swaths of underserved population. In my thesis, I hope to explore how to conduct bus network redesigns in a method that balances service quality and cost with equity and accessibility factors.

I hope to design a method for doing this and conduct a case study of a rural town, if possible, partnering with a transit agency that could benefit from this. I will begin by developing a novel framework for co-designing and modeling urban bus services, which offers the chance to reimagine how these systems are designed at the city-level scale. This framework should be able to ingest a variety of different networks, such as OpenStreetMap data, static GTFS feeds, and origin-destination pairs, into a common format. 
 
From there, using applied category theory for co-design, I can create a blueprint that will be able to ingest the network and capture the technical aspects of transit and the complex interplay of stakeholders, from city planners to operators to passengers. This formalism will help to think rigorously about the structure of urban mobility systems and how to optimize them, allowing us to approach and partner with a public transit agency (if possible, such as Metro McAllen) and apply this real-world model to assess feasibility and accuracy, as well as to identify areas for further improvement and research. We can also use this data to build test cases and stress-test our model under realistic urban scenarios.
